\documentclass[11pt]{article}
\usepackage{geometry}
\geometry{margin=1in}
\usepackage{amsmath,amssymb,amsthm}
\usepackage{hyperref}

\newtheorem{theorem}{Theorem}

\title{Disproof of conjecture \texttt{00000001011} \\ (finite unit-distance cliques)}
\author{earthking11}
\date{2026-10-01}

\begin{document}
\maketitle

\section{The conjecture}

The conjecture (file \texttt{conjectures/00000001011.md}) states:

\begin{quote}
\textbf{Conjecture.} In $\mathrm{PG}(2,q)$, the maximum clique of the
unit-distance graph (distance $\sqrt{d}$ for $d$ a quadratic residue of
$\mathbb{F}_q$) has size $q+1$ (finite unit-distance cliques).
\end{quote}

We disprove it already at $q=2$, where the maximum clique has size $2<3$.

\section{Modeling the distance}

The parenthetical definition in the conjecture is ambiguous, so we fix the
most natural reading.  Distances are defined on the affine patch
$A=\mathbb{F}_q^2$ of $\mathrm{PG}(2,q)$ via the standard quadratic form
\[
d(p,q)=(x_1-x_2)^2+(y_1-y_2)^2\in\mathbb{F}_q .
\]
``Distance $\sqrt{d}$ for $d$ a quadratic residue'' is read as: the squared
distance $d$ is a quadratic residue, so that the distance $\sqrt{d}$ exists
in $\mathbb{F}_q$.  We use the following model.

\begin{itemize}
\item \textbf{Model A (primary):} distinct points $p\ne q$ are adjacent
  iff $d(p,q)$ is a \emph{nonzero} quadratic residue of $\mathbb{F}_q$.
\item Points at infinity have no defined distance and are treated as
  isolated vertices.  A distance defined by a quadratic form on the affine
  patch cannot be extended to $\mathrm{PG}(2,q)$ in a projectively invariant
  way (scaling homogeneous coordinates changes the value of the form), so
  ``affine patch plus isolated points at infinity'' is the unique natural
  model.  Isolated vertices cannot raise the clique number, so the
  computation below is unaffected by this choice.
\end{itemize}

\section{The counterexample: $q=2$}

For $q=2$, the nonzero quadratic residues of $\mathbb{F}_2$ are
$\mathrm{QR}(\mathbb{F}_2)=\{1\}$.  The affine patch is
$\mathbb{F}_2^2=\{(0,0),(0,1),(1,0),(1,1)\}$.
Distinct points $p,q$ are adjacent iff
\[
d(p,q)\equiv (\Delta x)^2+(\Delta y)^2 \equiv 1 \pmod 2 .
\]
The six pairs give:
\begin{center}
\begin{tabular}{c|c|c}
pair & $d(p,q)\bmod 2$ & edge? \\ \hline
$(0,0)$--$(0,1)$ & $1$ & yes \\
$(0,0)$--$(1,0)$ & $1$ & yes \\
$(0,0)$--$(1,1)$ & $0$ & no \\
$(0,1)$--$(1,0)$ & $0$ & no \\
$(0,1)$--$(1,1)$ & $1$ & yes \\
$(1,0)$--$(1,1)$ & $1$ & yes
\end{tabular}
\end{center}
Hence the graph is the $4$-cycle
$(0,0)$--$(0,1)$--$(1,1)$--$(1,0)$--$(0,0)$:
it has $4$ edges and the two diagonals are non-edges.
The $\binom{4}{3}=4$ triples are therefore none of them cliques
(every triple contains a diagonal), so the maximum clique has size
$2<3=q+1$, contradicting the conjecture.  The full computation is
independently reproduced by the accompanying \texttt{reproduce.py}
(pure standard library) and formalized in the Lean~4 project
(\texttt{lean4/}).

\begin{theorem}
Under the model above, the unit-distance graph on the affine patch
$\mathbb{F}_2^2\subset\mathrm{PG}(2,2)$ has maximum clique size $2\ne 3$.
Hence the conjecture is false.
\end{theorem}

\section{Remarks}

\begin{enumerate}
\item The computation extends to $q=3,4,5,7,8,9$: under Model~A the maximum
  clique is exactly $q$ (the affine line $y=0$ is always a $q$-clique since
  $(t_1-t_2)^2$ is a nonzero square), never $q+1$.  The conjecture likely
  confuses an \emph{affine} line ($q$ points) with a \emph{projective} line
  ($q+1$ points).  We formalize only the $q=2$ instance, which suffices as a
  disproof; the values for larger $q$ are reported as computational
  observations, not theorems.
\item As a robustness check, we also tested the literal reading
  ``adjacent iff $d(p,q)=1$'' (Model~B): the maximum cliques for
  $q=2,3,4,5,7,8,9$ are $2,3,2,2,2,2,3$, again never $q+1$.  So the
  conjecture fails under both natural readings of the distance.
\end{enumerate}

\end{document}
